%% ============================================================
%%  Lecture 1 -- Tuesday, 26 May 2026
%%  Subfile of main.tex: compiles standalone OR as part of the book.
%%  Built from sources/2026-05-26/ (board-transcript.md, outline.md).
%% ============================================================
\documentclass[main.tex]{subfiles}

\begin{document}

\beginlecture{1}{Tuesday, 26 May 2026}%
  {From the Naturals to the Reals: Induction, Ordered Fields, and Completeness}

\epigraph{God made the integers; all else is the work of man.}{Leopold Kronecker}

\begin{lecturecontents}{Contents of this lecture}
  \item \hyperlink{1.1}{1.1\quad The natural numbers and induction}
  \item \hyperlink{1.2}{1.2\quad The integers and the rationals}
  \item \hyperlink{1.3}{1.3\quad Fields and ordered fields}
  \item \hyperlink{1.4}{1.4\quad Why the rationals are not enough}
  \item \hyperlink{1.5}{1.5\quad Bounds and the supremum}
  \item \hyperlink{1.6}{1.6\quad The real number system}
  \item \hyperlink{1.7}{1.7\quad The Archimedean property, density, and decimals}
\end{lecturecontents}

\bigskip
\noindent\textbf{Overview.} The course opens with a deceptively simple question---is
$0.999\ldots=1$?---whose honest answer requires building the real numbers from the ground up. We
begin with the natural numbers and induction, construct $\mathbb Z$ and $\mathbb Q$, and isolate
the field and order axioms they satisfy. The rationals then fall short for analysis: $\sqrt2$ is
missing, and the sets bounding it have neither a largest nor a smallest element. Repairing that
gap is the \emph{least-upper-bound property}, which characterises $\mathbb R$; from it follow the
Archimedean property and the density of $\mathbb Q$---and, at last, the resolution of
$0.999\ldots=1$.

\clearpage
% ============================================================
\sub{1.1}{The natural numbers and induction}

\begin{Obj}{Definition}{1.1.1}{The natural numbers}
The \emph{natural numbers} are $\mathbb N=\{0,1,2,3,\dots\}$\note{We take $0\in\mathbb N$; Rudin instead starts the positive integers at $1$. The convention is harmless---only the base of an induction shifts (Gauss's sum, \xref{1.1.4}, begins at $n=1$).}, generated from $0$ by the
\emph{successor function}
\[ S:\mathbb N\to\mathbb N,\qquad n\mapsto n+1 . \]
\end{Obj}

\begin{Obj}{Axiom}{1.1.2}{Induction}
If $A\subseteq\mathbb N$ satisfies $0\in A$ and $n\in A\Rightarrow n+1\in A$, then $A=\mathbb N$.
\end{Obj}

\begin{Obj}{Theorem}{1.1.3}{The principle of mathematical induction}
Let $P(n)$ be a statement for each $n\in\mathbb N$. If $P(0)$ holds, and $P(n)\Rightarrow P(n+1)$
for every $n$, then $P(n)$ holds for all $n\in\mathbb N$.
\end{Obj}
\begin{Prf}{1.1.3}{Proof of the Induction Principle from the Axiom.}{[Direct]\quad uses \xref{1.1.2}.}
\item Let $A=\{n\in\mathbb N : P(n)\text{ holds}\}$.
\item Since $P(0)$ holds, $0\in A$; and $P(n)\Rightarrow P(n+1)$ says $n\in A\Rightarrow n+1\in A$.
\item By the induction axiom (\xref{1.1.2}), $A=\mathbb N$; that is, $P(n)$ holds for every $n$.
\end{Prf}

\begin{ObjL}{Example}{1.1.4}{Gauss's sum}
For every integer $n\ge1$,\qquad $\displaystyle 1+2+\cdots+n=\frac{n(n+1)}{2}$.
\end{ObjL}
\begin{Prf}{1.1.4}{Proof by Induction of Gauss's Formula.}{[A10]\quad from the base $n=1$ (apply \xref{1.1.3} to the shifted statement).}
\item For $n=1$ the formula reads $1=\dfrac{1\cdot2}{2}$, which holds.
\item Assume $1+\cdots+n=\dfrac{n(n+1)}{2}$. Adding $n+1$ to both sides,
  \[ 1+\cdots+n+(n+1)=\frac{n(n+1)}{2}+(n+1)=\frac{(n+1)(n+2)}{2}, \]
  which is the formula for $n+1$. The result follows by induction (\xref{1.1.3}).
\end{Prf}

\begin{Fig}{1.1.5}{Induction as toppling dominoes}{The base case $P(0)$ topples the first domino; the inductive step $P(n)\Rightarrow P(n+1)$ passes the fall to the next, so every $P(n)$ falls in turn---the conclusion of \xref{1.1.2} and \xref{1.1.3}.}
\begin{tikzpicture}[x=1cm,y=1cm]
  \draw[cuBlueDk] (0.3,0) -- (7.4,0);
  \foreach \x/\n in {2/1,3/2,4/3,5/4,6/5}{
    \fill[cuBlueLt!70] (\x-0.09,0) rectangle (\x+0.09,0.8);
    \draw[cuBlueDk] (\x-0.09,0) rectangle (\x+0.09,0.8);
    \node[cuBlueDk,below=3pt] at (\x,0) {$\n$};
  }
  \begin{scope}[rotate around={-34:(1,0)}]
    \fill[figTerm] (0.91,0) rectangle (1.09,0.8);
    \draw[cuBlueDk] (0.91,0) rectangle (1.09,0.8);
  \end{scope}
  \node[figTerm,below=3pt] at (1,0) {$0$};
  \draw[->,figTerm,thick] (1.7,1.05) .. controls (3.3,1.5) and (4.9,1.5) .. (6.3,1.0);
  \node[figTerm] at (4,1.72) {$P(n)\Rightarrow P(n+1)$};
  \node[cuBlueDk] at (7.0,0.4) {$\cdots$};
\end{tikzpicture}
\end{Fig}

\begin{Fig}{1.1.6}{Gauss's sum as half a rectangle}{Two copies of the staircase $1+2+3+4$ interlock into a $4\times5$ grid, so $2(1+2+3+4)=4\cdot5$; in general $1+\cdots+n=\tfrac{n(n+1)}2$ (\xref{1.1.4}).}
\begin{tikzpicture}[x=0.6cm,y=0.55cm]
  \foreach \cx in {1,...,4}{
    \foreach \cy in {1,...,\cx}{\filldraw[figTerm] (\cx,\cy) circle (2.3pt);}
  }
  \foreach \cx in {1,...,4}{
    \pgfmathtruncatemacro{\st}{\cx+1}
    \foreach \cy in {\st,...,5}{\draw[figTerm,fill=cuBlueLt] (\cx,\cy) circle (2.3pt);}
  }
  \draw[cuBlueDk] (0.5,0.5) rectangle (4.5,5.5);
  \node[cuBlueDk,below=2pt] at (2.5,0.5) {$n=4$};
  \node[cuBlueDk,right] at (4.6,3) {$n+1=5$};
\end{tikzpicture}
\end{Fig}

% ============================================================
\sub{1.2}{The integers and the rationals}

\begin{Obj}{Construction}{1.2.1}{The integers}
The \emph{integers} are signed naturals with the two zeros glued together:
\[ \mathbb Z:=\big(\{+,-\}\times\mathbb N\big)\big/\!\sim,\qquad (+,0)\sim(-,0). \]
Writing $+n,-n$ for the classes recovers $\mathbb Z=\{\dots,-2,-1,0,1,2,\dots\}$.
\end{Obj}

\begin{Obj}{Construction}{1.2.2}{The rationals}
The \emph{rationals} are equivalence classes of integer pairs with nonzero second coordinate:
\[ \mathbb Q:=\{(p,q)\in\mathbb Z\times\mathbb Z : q\neq0\}\big/\!\sim,\qquad
   (p,q)\sim(r,s)\iff ps=rq . \]
The class of $(p,q)$ is written $\dfrac pq$.
\end{Obj}

\begin{ObjL}{Remark}{1.2.3}{One number, many fractions}
A rational \emph{is} the whole class, e.g.\
$\big[\tfrac12\big]=\big\{\tfrac12,\tfrac24,\tfrac{-1}{-2},\tfrac{17}{34},\dots\big\}$. In
practice one names it by its \emph{preferred representative}: the irreducible fraction
$\dfrac ab$ with $\gcd(a,b)=1$ and the sign carried in the numerator.
\end{ObjL}

% ============================================================
\sub{1.3}{Fields and ordered fields}

\begin{Obj}{Definition}{1.3.1}{Field}
A \emph{field} is a set $F$ with two operations $+:F\times F\to F$ and $\cdot:F\times F\to F$
satisfying eleven axioms\note{Five for addition (A1--A5), five for multiplication (M1--M5), and
the distributive law (D); the full statements are \R{1.12}.}---making $F$ an abelian group under
$+$, making $F\setminus\{0\}$ an abelian group under $\cdot$, with multiplication distributing
over addition.
\end{Obj}

\begin{Obj}{Definition}{1.3.2}{Order}
An \emph{order} on a set $S$ is a relation $<$ such that
\begin{enumerate}[label=\textup{(\arabic*)},leftmargin=2em,topsep=2pt,itemsep=2pt]
  \item (\emph{trichotomy}) for $x,y\in S$, exactly one of $x<y$, $y<x$, $x=y$ holds;
  \item (\emph{transitivity}) $x<y$ and $y<z$ imply $x<z$.
\end{enumerate}
A set carrying an order is \emph{totally} (or \emph{linearly}) ordered.
\end{Obj}

\begin{Obj}{Definition}{1.3.3}{Ordered field}
An \emph{ordered field} is a field $F$ with an order $<$ compatible with the operations:
\begin{enumerate}[label=\textup{(\arabic*)},leftmargin=2em,topsep=2pt,itemsep=2pt]
  \item $y<z\Rightarrow x+y<x+z$ for all $x$;
  \item $x>0$ and $y>0\Rightarrow xy>0$.
\end{enumerate}
\end{Obj}

\begin{ObjL}{Remark}{1.3.4}{The rationals are an ordered field}
$\mathbb Q$ is an ordered field. The rest of the lecture asks what it \emph{lacks}---and why that
makes it unfit for analysis.
\end{ObjL}

% ============================================================
\sub{1.4}{Why the rationals are not enough}

\begin{Obj}{Proposition}{1.4.1}{No rational squares to $2$}
There is no $x\in\mathbb Q$ with $x^2=2$.
\end{Obj}
\begin{Prf}{1.4.1}{Proof by Contradiction of the Irrationality of $\sqrt2$.}{[A9]\quad uses ``$n^2$ even $\Rightarrow n$ even''.}
\item Suppose $x=\dfrac pq\in\mathbb Q$ is irreducible with $x^2=2$. Then $p^2=2q^2$, so $p^2$ is
  even, hence $p$ is even (an odd integer has an odd square); write $p=2k$.
\item Substituting, $4k^2=2q^2$, so $q^2=2k^2$ is even, hence $q$ is even.
\item Then $2\mid\gcd(p,q)$, contradicting that $\dfrac pq$ is irreducible. So no such $x$ exists.
\end{Prf}

\begin{Obj}{Proposition}{1.4.2}{The rationals have a gap}
Let $A=\{p\in\mathbb Q : p>0,\ p^2<2\}$ and $B=\{p\in\mathbb Q : p>0,\ p^2>2\}$. Then $A$ has no
largest element and $B$ has no smallest element.
\end{Obj}
\begin{Prf}{1.4.2}{Proof by an Explicit Construction of the Gap.}{[Constr]\quad cf.\ \R{1.1}.}
\item Given $p>0$, set
  \[ q=p-\frac{p^2-2}{p+2}=\frac{2p+2}{p+2}>0 . \]
\item A computation gives
  \[ q^2-2=\frac{2\,(p^2-2)}{(p+2)^2}, \]
  so $q^2-2$ has the same sign as $p^2-2$: if $p\in A$ then $q\in A$, and if $p\in B$ then $q\in B$.
\item Also $q-p=-\dfrac{p^2-2}{p+2}$, so $q>p$ when $p\in A$ and $q<p$ when $p\in B$.
\item Hence every $p\in A$ is exceeded by a larger $q\in A$, so $A$ has no largest element; and
  every $p\in B$ is undercut by a smaller $q\in B$, so $B$ has no smallest element.
\end{Prf}

\begin{Fig}{1.4.3}{The gap at $\sqrt2$}{The positive rationals $A$ ($p^2<2$) climb toward $\sqrt2$ with no largest member; $B$ ($p^2>2$) descend toward it with no smallest; in $\mathbb Q$ the meeting point $\sqrt2$ is simply absent.}
\begin{tikzpicture}[x=1cm,y=1cm]
  \fill[cuBlueLt!65] (1.2,-0.12) rectangle (3.7,0.12);
  \fill[cuBlueLt!65] (4.3,-0.12) rectangle (6.6,0.12);
  \draw[->] (0.3,0) -- (7.3,0) node[right] {$\mathbb Q$};
  \filldraw (1,0) circle (1pt) node[below=4pt] {$0$};
  \node[cuBlueDk,above=3pt] at (2.45,0.12) {$A$};
  \node[cuBlueDk,above=3pt] at (5.45,0.12) {$B$};
  \draw[cuBlueDk,thick] (4,0) circle (2.4pt);
  \node[cuBlueDk,below=4pt] at (4,0) {$\sqrt2$ (absent)};
  \foreach \x in {2.0,2.9,3.4,3.66}{\filldraw[figTerm] (\x,0) circle (1pt);}
  \foreach \x in {4.34,4.6,5.1,6.0}{\filldraw[figTerm] (\x,0) circle (1pt);}
\end{tikzpicture}
\end{Fig}

% ============================================================
\sub{1.5}{Bounds and the supremum}

\begin{Obj}{Definition}{1.5.1}{Upper bound}
Let $S$ be an ordered set and $E\subseteq S$. An element $\beta\in S$ is an \emph{upper bound}
for $E$ if $x\le\beta$ for every $x\in E$.
\end{Obj}

\begin{Obj}{Definition}{1.5.2}{Least upper bound (supremum)}
$\beta$ is the \emph{least upper bound} of $E$ if it is an upper bound for $E$ and no smaller
element is: $\alpha<\beta$ implies $\alpha$ is not an upper bound. We write
$\beta=\sup E=\operatorname{lub}E$.
\end{Obj}

\begin{Obj}{Definition}{1.5.3}{Lower bound and infimum}
Dually, $\beta\in S$ is a \emph{lower bound} for $E$ if $\beta\le x$ for every $x\in E$, and the
\emph{greatest lower bound} (or \emph{infimum}) if no larger element is a lower bound: $\beta<\alpha$
implies $\alpha$ is not a lower bound. We write $\beta=\inf E=\operatorname{glb}E$. Since
$\inf E=-\sup(-E)$, the least-upper-bound property (\xref{1.6.1}) carries a matching
greatest-lower-bound property.
\end{Obj}

% ============================================================
\sub{1.6}{The real number system}

\begin{Obj}{Theorem}{1.6.1}{The real number system}
There is an ordered field $\mathbb R\supseteq\mathbb Q$ in which every nonempty subset that is
bounded above has a least upper bound. It is unique up to isomorphism.\note{An isomorphism of ordered fields is a bijection preserving $+$, $\cdot$, and $<$; ``unique up to isomorphism'' means any two complete ordered fields are indistinguishable as such. Two were built in $1872$: Dedekind's cuts (\emph{Stetigkeit und irrationale Zahlen}) and Cantor's equivalence classes of Cauchy sequences of rationals.}
\end{Obj}

\begin{ObjL}{Remark}{1.6.2}{Completeness and density}
The property in \xref{1.6.1} is \emph{completeness}: $\mathbb R$ is complete iff it has the
least-upper-bound property. A first consequence is that $\mathbb Q$ is \emph{dense} in
$\mathbb R$---between any two reals lies a rational (\xref{1.7.2}).
\end{ObjL}

\begin{ObjL}{Remark}{1.6.3}{Three ways to build $\mathbb R$}
\begin{itemize}[leftmargin=1.3em,topsep=2pt,itemsep=1pt,label=\textbullet]
  \item \emph{Accept it}---take \xref{1.6.1} on faith and learn to use $\sup$.
  \item \emph{Dedekind cuts} (\xref{1.6.4})---realise each real as a downward-closed set of rationals.
  \item \emph{Completion}---the general procedure that completes any metric space (deferred).
\end{itemize}
\end{ObjL}

\begin{Obj}{Definition}{1.6.4}{Dedekind cut}
A \emph{cut} is a set $\alpha\subseteq\mathbb Q$ such that
\begin{enumerate}[label=\textup{(\arabic*)},leftmargin=2em,topsep=2pt,itemsep=2pt]
  \item $\alpha\neq\varnothing$ and $\alpha\neq\mathbb Q$;
  \item $\alpha$ is downward closed: $r\in\alpha$ and $s<r$ imply $s\in\alpha$;
  \item $\alpha$ has no largest element.
\end{enumerate}
The cuts, suitably ordered, \emph{are} the real numbers.
\end{Obj}

\begin{ObjL}{Example}{1.6.5}{$\sqrt2$ as a cut}
The real number $\sqrt2$ is the cut
\[ \alpha=\{q\in\mathbb Q : q<0\}\ \cup\ \{q\in\mathbb Q : q^2<2\} . \]
This is the set $A$ of \xref{1.4.2} together with the non-positive rationals; the former ``gap''
is now a bona fide element.
\end{ObjL}

% ============================================================
\sub{1.7}{The Archimedean property, density, and decimals}

\begin{Obj}{Theorem}{1.7.1}{Archimedean property}
If $x,y\in\mathbb R$ and $x>0$, then there is a positive integer $n$ with $nx>y$.\note{Named by Otto Stolz in the $1880$s after Archimedes, who postulated it in \emph{On the Sphere and the Cylinder}; the idea is older still---Euclid records it (\emph{Elements} V, Def.~4) and credits Eudoxus.}
\end{Obj}
\begin{Prf}{1.7.1}{Proof by Contradiction of the Archimedean Property.}{[A5$+$A9]\quad via the supremum; uses \xref{1.6.1}.}
\item Suppose not: the set $A=\{nx : n\in\mathbb N\}$ is bounded above. By the least-upper-bound
  property (\xref{1.6.1}) it has a supremum $\alpha=\sup A$.
\item Since $x>0$, $\alpha-x<\alpha$, so $\alpha-x$ is not an upper bound of $A$: some
  $m\in\mathbb N$ has $mx>\alpha-x$.
\item Then $(m+1)x>\alpha$, yet $(m+1)x\in A$---contradicting $\alpha=\sup A$. So $A$ is not
  bounded above, and in particular some $nx>y$.
\end{Prf}

\begin{Obj}{Theorem}{1.7.2}{Density of $\mathbb Q$ in $\mathbb R$}
If $x,y\in\mathbb R$ and $x<y$, then there is a rational $q$ with $x<q<y$.
\end{Obj}
\begin{Prf}{1.7.2}{Proof by the Archimedean Property of the Density of $\mathbb Q$.}{[Direct]\quad uses \xref{1.7.1}.}
\item Since $y-x>0$, the Archimedean property (\xref{1.7.1}) gives $n\in\mathbb N$ with
  $n(y-x)>1$, i.e.\ $ny-nx>1$.
\item Since $ny-nx>1$, the interval $(nx,ny)$ contains an integer $m$ (\appref{1.A.1}).
\item Dividing by $n>0$, $\ x<\dfrac mn<y$, so $q=\dfrac mn$ works.
\end{Prf}

\begin{Obj}{Construction}{1.7.3}{Real numbers as decimals}
Every \emph{nonnegative} real is the supremum of its finite decimal truncations (a negative real
is then named by its sign):
\[ 1.4112\ldots\ \longmapsto\ \sup\{\,1,\ 1.4,\ 1.41,\ 1.411,\ \dots\,\}. \]
\end{Obj}

\begin{ObjL}{Example}{1.7.4}{$0.999\ldots=1$}
This settles the opening question. Put $x=1-0.999\ldots\,$; then $0\le x<10^{-k}$ for every
$k\in\mathbb N$. By the Archimedean property (\xref{1.7.1})\note{If $x>0$, the Archimedean property gives $n$ with $nx>1$; taking $k$ with $10^k\ge n$ then yields $10^k x>1$, i.e.\ $x>10^{-k}$.} no positive real lies below
every $10^{-k}$, so $x=0$---that is, $0.999\ldots=1$. More generally every terminating decimal equals
its repeating-$9$ twin ($2.5=2.4999\ldots$, and so on), so
\[ \mathbb R=\{\text{decimal expansions}\}\big/\!\sim, \]
the quotient that identifies each terminating decimal with that twin.
\end{ObjL}

\begin{Fig}{1.7.5}{Density of $\mathbb Q$ in $\mathbb R$}{Choosing $n$ with $1/n<y-x$ makes the grid of multiples $m/n$ finer than the gap, so some $m/n$ lands strictly between $x$ and $y$ (\xref{1.7.2}).}
\begin{tikzpicture}[x=1cm,y=1cm]
  \fill[figLim!22] (4.1,-0.1) rectangle (5.5,0.1);
  \draw[->] (0.3,0) -- (8.0,0) node[right]{$\mathbb R$};
  \foreach \x in {0.7,1.4,2.1,2.8,3.5,4.2,4.9,5.6,6.3,7.0,7.7}{
    \draw[cuBlueDk!55] (\x,0.13) -- (\x,-0.13);
  }
  \node[cuBlueDk!75,below=9pt] at (1.75,0) {grid $\tfrac{m}{n}$, spacing $\tfrac1n$};
  \filldraw (4.1,0) circle (1.1pt) node[black,above=3pt]{$x$};
  \filldraw (5.5,0) circle (1.1pt) node[black,above=3pt]{$y$};
  \filldraw[figTerm] (4.9,0) circle (1.7pt) node[figTerm,below=4pt]{$q=\tfrac{m}{n}$};
\end{tikzpicture}
\end{Fig}

% ============================================================
\lectureappendix

\begin{ObjL}{Supplement}{1.A.1}{An interval longer than $1$ contains an integer}
If $s,t\in\mathbb R$ with $t-s>1$, then there is an integer $m$ with $s<m<t$. (Used silently in
the density proof, \xref{1.7.2}.)
\end{ObjL}
\begin{Prf}{1.A.1}{Direct Proof from the Archimedean Property.}{[Direct]\quad uses \xref{1.7.1}; cf.\ \R{1.20}.}
\item The integers exceeding $s$ are nonempty (by the Archimedean property) and bounded below by
  $s$, so they have a least element $m$; thus $m-1\le s<m$.
\item Hence $m\le s+1<t$, and therefore $s<m<t$.
\end{Prf}

\begin{ObjL}{Supplement}{1.A.2}{Every nonnegative real has a decimal expansion}
Construction \xref{1.7.3} names each nonnegative real as the supremum of its truncations; here we
build those truncations digit by digit and confirm the supremum is the number we began with. Given
$x\ge 0$, there are an integer $a_0\ge 0$ and digits $d_1,d_2,\dots\in\{0,\dots,9\}$ whose
truncations
\[ t_k\;=\;a_0+\sum_{i=1}^{k}\frac{d_i}{10^{i}} \qquad\text{satisfy}\qquad 0\le x-t_k<10^{-k}, \]
so that $x=\sup_k t_k$. The rule below always takes the \emph{largest} admissible digit, so a
number with a terminating expansion receives it ($1=1.000\ldots$), and the repeating-$9$ twin
(\xref{1.7.4}) is the single expansion this rule never produces.
\end{ObjL}
\begin{Prf}{1.A.2}{Direct Proof by Greedy Digit Selection.}{[Direct]\quad uses \xref{1.7.1}, \xref{1.A.1}.}
\item For the integer part, let $a_0$ be the greatest integer with $a_0\le x$ --- it exists by the
  least-integer argument of \xref{1.A.1}, applied to the integers exceeding $x$, and $a_0\ge 0$
  since $x\ge 0$. Then $t_0=a_0$ obeys $0\le x-t_0<1$.
\item For the digits, suppose $t_{k-1}$ has been built with $0\le x-t_{k-1}<10^{-(k-1)}$. Let $d_k$ be
  the largest digit in $\{0,\dots,9\}$ with $t_{k-1}+d_k\,10^{-k}\le x$; the digit $0$ qualifies, and
  $d_k\le 9$ because $x<t_{k-1}+10\cdot 10^{-k}=t_{k-1}+10^{-(k-1)}$. Put $t_k=t_{k-1}+d_k\,10^{-k}$.
  Then $t_k\le x$, while maximality of $d_k$ gives $x<t_k+10^{-k}$; hence $0\le x-t_k<10^{-k}$ and
  the invariant persists.
\item For the supremum: every $t_k\le x$, so $x$ is an upper bound. If $\beta<x$, then
  $x-\beta>10^{-k}$ for some $k$ (Archimedean property, \xref{1.7.1}), so
  $t_k>x-10^{-k}>\beta$; thus no $\beta<x$ bounds the truncations, and $x=\sup_k t_k$.
\end{Prf}

% per-lecture Index of Objects -- only when compiling this lecture on its own;
% in the book the master index in main.tex covers all lectures.
\ifSubfilesClassLoaded{\clearpage\objindex}{}

\end{document}
